
\begin{algorithm}[t]
    
    \caption{\texttt{PrE-Text} \label{private evolution}}
    \begin{algorithmic}[1]
    \Statex {\bfseries Input:}
     Number of iterations: $T$\\
     \hspace{\algorithmicindent} Number of generated samples: $N_{\text{syn}}$\\
     \hspace{\algorithmicindent} Initial population: $S_1$ \label{line:initialpop} \\
     {\bfseries Output:} Synthetic data $S_{\text{syn}}$ 
    \For{$t \gets 1 \dots T$} 
        \State \textsc{// see \cref{algo:dphistogram}}
        \State $\text{hist}_t \gets \texttt{Histogram}(S_t)$ \label{line:dp_hist}
        \State $P_t$ $\gets$ $\text{hist}_t / \text{sum}(\text{hist}_t)$
        \State $S_t'$ $\gets$ draw $N_{\text{syn}}$ samples from $P_t$ \label{line:surviving}
        \State \algemph{$S_t \gets \texttt{Variation}(S_t')$ }\label{line:variation}
    \EndFor

      \State  \algemph{$S_{\text{syn}} \gets \texttt{Expand}(\cup_{t=1}^T \texttt{Set}(S_t'))$ \label{line:expand}}
    \State \Return $S_{\text{syn}}$
    \end{algorithmic}
\end{algorithm}

\begin{algorithm}[t]
    \caption{\texttt{Histogram}}
    \begin{algorithmic}[1]
    \State {\bfseries Settings:}  Private samples: Clients $\{C_i\}_{i \in \mathcal{C}_{\text{train}}}$\\
     \hspace{\algorithmicindent} Noise multiplier for histogram: $\sigma$\\
     \hspace{\algorithmicindent} Number of generated samples: $N_{\text{syn}}$\\
     \hspace{\algorithmicindent} Threshold for histogram: $H$ \\
     \hspace{\algorithmicindent} Distance function: $d(\cdot, \cdot)$ 
    \State {\bfseries Input:} Generated samples $S = \{z_i\}_{i=1}^{N_{\text{syn}}}$ \\
     {\bfseries Output:} Nearest neighbors histogram on $S$
    \For{$i \in \mathcal{C}_{\text{train}}$}
        \State {$\text{hist}_i \gets [0, \dots, 0]$ \label{line:fitness}}
        \For{$x_{\text{priv}} \in \mathcal{D}_i$}
            \State $l \gets \text{argmin}_{j \in [N_{\text{syn}}]} d(x_{\text{priv}}, z_j)$ \label{line:dpnn}
            \State $\text{hist}_i[l] \gets \text{hist}_i[l] + 1$
        \EndFor
        \State $\text{hist}_i \gets \text{hist}_i + \mathcal{N}(0, (\sigma^2/|\mathcal{C}_{\text{train}}|) I_{N_{\text{syn}}})$ \label{line:dp}
    \EndFor
    \State $\text{hist} \gets \sum_{i \in \mathcal{C}_{\text{train}}} \text{hist}_i$ \label{line:collect} 
    \State $\text{hist} \gets \text{max}(\text{hist} - H, 0)$ \label{line:threshold} (elementwise subtraction)
    \State \Return \text{hist}
    \end{algorithmic}
    \label{algo:dphistogram}
\end{algorithm}
\section{PrE-Text}
\label{sec: alg description}
The main intuition of PrE-Text (and PE) is that public foundation models should be capable of producing samples that are similar to the private data with some non-negligible probability.  Therefore, to generate DP synthetic data similar to private data, we steer the foundation model (privately) towards the user data in a multi-round process.  Briefly, we make several important changes to the PE algorithm: (1) we adapt it from the image setting to the text setting; (2) we exploit synthetic data from the \emph{entire} PE process, rather than only the last round; and (3) we add a post-processing phase that utilizes the output of PE as seeds for another LLM. We provide pseudocode in \cref{private evolution}, and highlight the new contributed steps in color. We now explain \cref{private evolution}.

\textbf{(1) Population of samples.} We start with an initial population of samples $S_1$ (\cref{line:initialpop}).  These samples can come from many different sources as long as they do not contain private information: for example public samples collected from the internet or samples randomly generated by a public generative foundation model.

\textbf{(2) Clients vote for the best synthetic samples.}  For round $t\geq 1$, we determine which of the generated samples $S_t$ represent the private samples the best.  We send all the generated samples $S_t$ to each client.  Each client counts for each generated sample $s\in S_t$ how many private samples had $s$ as their nearest neighbor in $S_t$.
The higher this count is, the ``better'' a generated sample is.  Thus, each client produces a nearest neighbors histogram with $|S_t|$ entries (\cref{line:fitness}). We determine nearest neighbor according to a distance function $d(y,z) = \|\Phi(y) - \Phi(z)\|_2$, where $\Phi$ is an embedding model.  


\textit{Lookahead.} To more accurately assess the closeness of a synthetic sample to a private sample, we amend the distance function from $d(y,z) = \|\Phi(y) - \Phi(z)\|_2$ to $d(y, z) = \|\Phi(y) - \frac{1}{K}\sum_{i=1}^K \Phi(z^i)\|_2$, where $z^0, \dots, z^K$ are $K$ variations of $z$ produced by using \texttt{Variation}. \citet{lin2023differentially} also uses this modification. Instead of sending the actual generated samples directly to all the clients, we send $\frac{1}{K}\sum_{i=1}^K \Phi(z^i)$ to all the clients for every $i$ for their nearest neighbors calculation \pcref{line:dpnn}.

\textit{(a) DP Noise.} In \cref{line:dp} each client adds noise to their nearest neighbors histogram to ensure DP. We compute client-level sensitivity assuming a known upper bound on the number of samples per client (e.g., via thresholding). 

\textit{(b) Federated Secure Aggregation.}  In \cref{line:collect}  we securely aggregate the histograms across the users.  Because the generated samples given to all the clients are the same, we sum the histograms using secure aggregation \citep{secagg} to get an aggregate histogram. 
% representing the ``best'' generated samples.

\textit{(c) Thresholding.} When we generate many samples, the majority of the probability mass of the histogram will be noise.  We improve the signal-to-noise ratio by thresholding the histogram at $H$ in \cref{line:threshold} \citep{lin2023differentially}. 

\textbf{(3) Use votes to choose the surviving samples.} We sample from the nearest neighbors histogram to produce surviving generated samples in \cref{line:surviving}, $S_{t+1}'$.  This new list of generated samples (there may be duplicates) tends to give more representation to generated samples that had more private samples close to them.

\textbf{(4) Produce variations of surviving samples.} We use \texttt{Variation} to generate a variation of the surviving samples $S_{t+1}'$ as the new population of samples (\cref{line:variation}). In PE, this was accomplished using diffusion models, which cannot be used for text.
We instead implement \texttt{Variation}  as follows. For each sample in $S'_t$, we produce a variation of it by masking $\textsc{mask}$\% of the tokens randomly, then filling in those masked tokens in with a masked language model. The resulting sample is then masked and filled-in again. This mask-fill process happens $W_{\text{steps}}$ times before returning the variation. We use RoBERTa-large \citep{roberta} as the masked language model.

\textbf{(5) Making efficient use of iterates.} We make several major modifications to the core Private Evolution algorithm to improve our usage of the iterates. First, \citet{lin2023differentially} use the final $S_T$ (\cref{line:variation}) as the synthetic dataset.  However, $S_T'$ contains more information about the private dataset than $S_T$ because \texttt{Variation} destroys information.  Therefore, we use $S_T'$ instead of $S_T$. Second, we find that $S_t'$ for $t=1,\dots,T$ all have valuable synthetic samples and show significant diversity between iterations.  Therefore, we choose to utilize $\bigcup_{t=1}^T \texttt{Set}(S_t')$ \pcref{line:expand} instead of just $S_T'$.  This more effectively utilizes the privacy budget.

\textbf{(6) Post-processing: Use LLM to expand the DP seed set.} 
PE originally used the final $S_T$ as the target synthetic data. 
However, we find that these samples are not high enough quality to fine-tune an LLM. 
We instead use \texttt{Expand} to generate more samples similar to the synthetic samples $\bigcup_{t=1}^T \texttt{Set}(S_t')$, which we use as our (DP) seed set to prompt the LLM.  \texttt{Expand} utilizes the synthetic data generation capabilities of large language models (LLMs) to generate a larger and more useful synthetic dataset.  Note that by the post-processing property of differential privacy \citep{dwork2006differential}, \texttt{Expand} will not leak any additional privacy.  Next we describe how \texttt{Expand} works.

Inspired by highly successful synthetic text generation \citep{alpaca, wang2022self, honovich2022unnatural, roziere2023code}, we generate synthetic text by using large foundational language models.  We follow a process similar to \citet{honovich2022unnatural}: for each synthetic sample to generate, we randomly choose three text samples to emulate from $\bigcup_{t=1}^T \texttt{Set}(S_t')$ (our DP seed set), and ask the LLM to generate a similar sample.  We use open-source LLaMA-2-7B \citep{touvron2023llama} as our large language model.  We provide the full prompt \pcref{fig:prompt}.

\subsection{Privacy analysis}
As noted in \citet{lin2023differentially}, the only function that utilizes private information is \cref{algo:dphistogram}. The DP histogram (\cref{line:collect}) contains private information.  As in \citet{lin2023differentially}, we use the Gaussian mechanism with constant noise multiplier $\sigma$ each time we receive a histogram. Since this is a Gaussian mechanism, we can use the moments accountant from the Opacus library \citep{opacus}. 
Details  can be found in \cref{app:privacy}.

% In particular, we use \textsc{opacus.accountants.analysis.rdp} to compute the privacy guarantee for PrE-Text (we input $T=11$ steps, $q=1.0$, and set the \textsc{noise\_multiplier} to be the ratio of $\sigma$ to the sensitivity (the max number of samples per client for PrE-Text), setting $\sigma$ to the value that gets us the desired $\epsilon$ value).
